\documentclass[11pt,a4paper]{article}

% ---------------------------------------------------------------------------
% Pakete
% ---------------------------------------------------------------------------
\usepackage[utf8]{inputenc}
\usepackage[T1]{fontenc}
\usepackage[a4paper,margin=2.4cm]{geometry}
\usepackage{lmodern}
\usepackage{microtype}
\usepackage{amsmath}
\usepackage{amssymb}
\usepackage{booktabs}
\usepackage{array}
\usepackage{xcolor}
\usepackage{listings}
\usepackage{enumitem}
\usepackage{tikz}
\usepackage{pgfplots}
\pgfplotsset{compat=1.18}
\usepackage[hidelinks]{hyperref}
\usepackage{fancyhdr}

\usetikzlibrary{arrows.meta,positioning,shapes.geometry,fit,backgrounds,calc}

% ---------------------------------------------------------------------------
% Farben & Stil
% ---------------------------------------------------------------------------
\definecolor{accent}{HTML}{2C6E9B}
\definecolor{accent2}{HTML}{3FA796}
\definecolor{accent3}{HTML}{E08E45}
\definecolor{ink}{HTML}{1F2933}
\definecolor{soft}{HTML}{F0F4F8}
\definecolor{codebg}{HTML}{F7F8FA}
\definecolor{codekw}{HTML}{2C6E9B}
\definecolor{codestr}{HTML}{3FA796}
\definecolor{codecom}{HTML}{8895A7}

\lstdefinestyle{py}{
  language=Python,
  basicstyle=\ttfamily\footnotesize\color{ink},
  keywordstyle=\color{codekw}\bfseries,
  stringstyle=\color{codestr},
  commentstyle=\color{codecom}\itshape,
  backgroundcolor=\color{codebg},
  frame=single,
  rulecolor=\color{soft},
  numbers=left,
  numberstyle=\tiny\color{codecom},
  showstringspaces=false,
  breaklines=true,
  columns=fullflexible,
  keepspaces=true,
  tabsize=4,
}

\pagestyle{fancy}
\fancyhf{}
\lhead{\small\color{accent}RAG Backend}
\rhead{\small\color{accent}AP 1A \textendash{} ChromaDB Integration \& Retrieval}
\cfoot{\thepage}
\renewcommand{\headrulewidth}{0.4pt}

% TikZ Knoten-Stile
\tikzset{
  layer/.style={draw=accent, fill=soft, rounded corners=3pt, line width=0.7pt,
                inner sep=8pt, align=center},
  comp/.style={draw=accent, fill=white, rounded corners=2pt, line width=0.6pt,
               minimum height=8mm, minimum width=26mm, align=center,
               font=\small},
  ext/.style={draw=accent3, fill=accent3!10, rounded corners=2pt, line width=0.6pt,
              minimum height=8mm, minimum width=26mm, align=center, font=\small},
  store/.style={draw=accent2, fill=accent2!12, cylinder, shape border rotate=90,
                aspect=0.25, minimum height=12mm, minimum width=22mm, align=center,
                font=\small, line width=0.6pt},
  flow/.style={-{Stealth[length=2.2mm]}, line width=0.8pt, draw=ink},
  flowd/.style={-{Stealth[length=2.2mm]}, line width=0.8pt, draw=accent3, dashed},
  lbl/.style={font=\scriptsize\itshape, color=codecom},
}

% ---------------------------------------------------------------------------
\begin{document}

\begin{titlepage}
\centering
\vspace*{2.2cm}
{\color{accent}\rule{\linewidth}{1.4pt}}\\[0.5cm]
{\Huge\bfseries RAG Backend\par}
\vspace{0.35cm}
{\LARGE Arbeitspaket 1A\par}
\vspace{0.2cm}
{\Large ChromaDB Integration \& Retrieval Validation\par}
{\color{accent}\rule{\linewidth}{1.4pt}}\\[1.1cm]

\begin{minipage}{0.85\linewidth}
\centering\large
Dokumentation der Implementierung des Dokument-Ingestion- und Vektor-Retrieval-Fundaments
für das on-premises Retrieval-Augmented-Generation-System.
\end{minipage}

\vfill
\begin{tabular}{rl}
\toprule
GitHub Issue & \#19 (Spur A, blocked by AP 0)\\
Modul-Scope & \texttt{backend/}\\
Embedding-Modell & \texttt{all-MiniLM-L6-v2} (384 Dim.)\\
Vektor-Store & ChromaDB (Cosine, persistent)\\
Top-1-Hit-Rate & \textbf{100\,\% (20/20)}\\
\bottomrule
\end{tabular}

\vspace{1.2cm}
{\color{codecom}\today}
\end{titlepage}

\tableofcontents
\newpage

% ===========================================================================
\section{Überblick \& Zielsetzung}
% ===========================================================================
Arbeitspaket~1A legt das \emph{Infrastruktur-Fundament} der Spur~A des RAG-Backends:
das Einlesen von Dokumenten, ihre Zerlegung in Chunks, deren Vektorisierung und die
persistente Ablage in einem Vektor-Store mit Cosine-Similarity. Es ist die
Voraussetzung für die nachgelagerte Streaming-Chat-Integration (AP~2A) und die
vollständige RAG-Chain (AP~4).

\paragraph{Definition of Done.} Alle Kriterien aus dem Implementierungsplan sind erfüllt:

\begin{center}
\renewcommand{\arraystretch}{1.25}
\begin{tabular}{>{\raggedright\arraybackslash}p{8.8cm} c}
\toprule
\textbf{Kriterium} & \textbf{Status}\\
\midrule
20-Dokument Ground-Truth-Korpus in ChromaDB ladbar & \checkmark\ 20/20\\
Top-1-Hit-Rate $\geq 80\,\%$ auf Ground-Truth-Queries & \checkmark\ \textbf{100\,\%}\\
Cosine-Similarity-Score $\geq$ Threshold (0.25) & \checkmark\ min 0.48 / $\varnothing$ 0.66\\
Unit-Tests für loader, chunker, embedder mit Mock-ChromaDB & \checkmark\ 30 Tests\\
\texttt{flake8}, \texttt{mypy}, \texttt{bandit} clean & \checkmark\\
Coverage $\geq 80\,\%$ & \checkmark\ 90\,--\,100\,\%\\
\bottomrule
\end{tabular}
\end{center}

% ===========================================================================
\section{Architektur}
% ===========================================================================
Die Implementierung folgt dem etablierten Schichtenmuster des Backends
(\texttt{services/core}, \texttt{services/dependency}, \texttt{services/utils}).
Fachlogik (Ingestion) ist von der Infrastruktur (Vektor-DB-Client) getrennt; alle
externen Abhängigkeiten werden injiziert, um Testbarkeit ohne laufenden Vektor-Store
und ohne Modell-Download zu garantieren.

\subsection{Schichten- und Komponentenmodell}

\begin{center}
\begin{tikzpicture}[node distance=6mm]

% --- core / ingestion layer ---
\node[comp] (loader)  {\textbf{loader.py}\\\scriptsize PDF/DOCX/TXT $\to$ Text};
\node[comp, right=of loader] (chunker) {\textbf{chunker.py}\\\scriptsize Recursive Split 512/50};
\node[comp, right=of chunker] (embedder) {\textbf{embedder.py}\\\scriptsize Embed + Upsert};

\begin{scope}[on background layer]
\node[layer, fit=(loader)(chunker)(embedder),
      label={[font=\small\bfseries,color=accent]above:services/core/ingestion}] (core) {};
\end{scope}

% --- dependency layer ---
\node[comp, below=16mm of chunker] (vectordb) {\textbf{vectordb.py}\\\scriptsize VectorDBClient (Singleton)};
\begin{scope}[on background layer]
\node[layer, fit=(vectordb),
      label={[font=\small\bfseries,color=accent]above:services/dependency}] (dep) {};
\end{scope}

% --- external ---
\node[ext, left=22mm of vectordb] (st) {\textbf{sentence-\\transformers}\\\scriptsize all-MiniLM-L6-v2};
\node[store, right=22mm of vectordb] (chroma) {ChromaDB\\\scriptsize Cosine\\\scriptsize persistent};

% edges
\draw[flow] (loader) -- (chunker);
\draw[flow] (chunker) -- (embedder);
\draw[flow] (embedder.south) -- ++(0,-3mm) -| (vectordb.north);
\draw[flowd] (embedder.west) to[bend right=12] (st.north);
\draw[flow] (vectordb) -- (chroma);
\draw[flowd] (st.south) |- (vectordb.west);

\node[lbl, below=1mm of dep] {Dependency-Injection: Client \& Modell sind austauschbar (real / fake)};
\end{tikzpicture}
\end{center}

\noindent\textcolor{accent3}{\rule{6mm}{2pt}}~\footnotesize externe Plattform-Komponente\quad
\textcolor{accent}{\rule{6mm}{2pt}}~Custom-Komponente\quad
\textcolor{accent2}{\rule{6mm}{2pt}}~Datenspeicher
\normalsize

\subsection{Ingestion-Datenfluss}
Der Schreibpfad transformiert eine Datei in durchsuchbare Vektoren. Jeder Schritt ist
eine reine, isoliert testbare Funktion bzw. Klasse.

\begin{center}
\begin{tikzpicture}[node distance=9mm and 12mm]
\node[ext] (file) {Datei\\\scriptsize .pdf/.docx/.txt};
\node[comp, right=of file] (text) {Reiner Text\\\scriptsize \texttt{extract\_text()}};
\node[comp, right=of text] (chunks) {Chunks\\\scriptsize \texttt{Chunk(source,idx,text)}};
\node[comp, right=of chunks] (vec) {Embeddings\\\scriptsize $384$-dim, normalisiert};
\node[store, right=of vec] (db) {ChromaDB\\\scriptsize upsert};

\draw[flow] (file) -- node[lbl,above]{loader} (text);
\draw[flow] (text) -- node[lbl,above]{chunker} (chunks);
\draw[flow] (chunks) -- node[lbl,above]{embedder} (vec);
\draw[flow] (vec) -- node[lbl,above]{id+meta} (db);
\end{tikzpicture}
\end{center}

\noindent Die Chunk-ID ist deterministisch und inhaltsadressiert
(\texttt{source:index:sha256[:16]}); ein erneutes Hochladen identischer Inhalte führt
damit zu idempotenten Upserts statt Duplikaten.

\subsection{Retrieval-Fluss}
Der Lesepfad bettet die Anfrage mit demselben Modell ein und ermittelt über
Cosine-Distance die ähnlichsten Chunks.

\begin{center}
\begin{tikzpicture}[node distance=8mm and 13mm]
\node[ext] (q) {Query};
\node[comp, right=of q] (qe) {Query-Embedding\\\scriptsize \texttt{embed\_query()}};
\node[store, right=of qe] (db) {ChromaDB\\\scriptsize \texttt{collection.query}};
\node[comp, right=of db] (res) {Top-$k$ Chunks\\\scriptsize +\,Metadaten,\,Distanz};

\draw[flow] (q) -- (qe);
\draw[flow] (qe) -- node[lbl,above]{$\vec{q}$} (db);
\draw[flow] (db) -- node[lbl,above]{cosine}(res);
\end{tikzpicture}
\end{center}

\noindent Ähnlichkeit $= 1 - \text{Distanz}$. Da alle Vektoren $\ell_2$-normalisiert sind,
entspricht das dem Skalarprodukt $\langle \vec{q}, \vec{d}\rangle$.

% ===========================================================================
\section{Komponenten im Detail}
% ===========================================================================

\subsection{\texttt{vectordb.py} \textendash{} Vektor-DB-Client}
Dünner Singleton-Wrapper um ChromaDBs \texttt{PersistentClient}. Die Collection wird
mit Cosine-Distanz angelegt; der Client ist injizierbar, der Persistenzpfad kommt aus
\texttt{CHROMA\_PATH} (Fallback \texttt{./chroma\_db}).

\begin{lstlisting}[style=py]
def get_collection(self) -> Collection:
    return self._client.get_or_create_collection(
        name=self._config.collection_name,
        metadata={"hnsw:space": "cosine"},
    )
\end{lstlisting}

\subsection{\texttt{loader.py} \textendash{} Dokument-Loader}
Dispatch nach Datei-Endung; pro Format ein lazy importierter Parser
(\texttt{pypdf}, \texttt{python-docx}, UTF-8). Nicht unterstützte Typen werfen
\texttt{UnsupportedDocumentError}. Eine bytes-basierte API (\texttt{extract\_text})
entkoppelt das Parsen vom Dateisystem und erleichtert den späteren Upload-Endpunkt.

\subsection{\texttt{chunker.py} \textendash{} Text-Splitter}
Kapselt LangChains \texttt{RecursiveCharacterTextSplitter} mit
\texttt{chunk\_size=512}, \texttt{chunk\_overlap=50}. Validiert die Konfiguration und
liefert indizierte \texttt{Chunk}-Records, die ihren Quellnamen tragen.

\subsection{\texttt{embedder.py} \textendash{} Embedding \& Upsert}
Zwei Verantwortlichkeiten, sauber getrennt: \texttt{EmbeddingService} lädt das
sentence-transformers-Modell \emph{lazy} (Tests injizieren ein deterministisches
Fake-Modell); \texttt{DocumentEmbedder} orchestriert Embedding und Upsert in eine
\emph{injizierte} Collection.

\begin{lstlisting}[style=py]
ids = [build_chunk_id(chunk) for chunk in chunks]
metadatas = [{"source": c.source, "chunk_index": c.chunk_index} for c in chunks]
embeddings = self._embeddings.embed_texts([c.text for c in chunks])
self._collection.upsert(ids=ids, embeddings=embeddings,
                        documents=documents, metadatas=metadatas)
\end{lstlisting}

% ===========================================================================
\section{Teststrategie}
% ===========================================================================
Zwei Test-Ebenen trennen schnelle, deterministische Logik-Tests von der echten,
schwergewichtigen Modell-Validierung.

\begin{center}
\begin{tikzpicture}[node distance=7mm and 16mm]
% mocked
\node[comp] (ut) {Unit-Tests\\\scriptsize 30 Tests};
\node[ext, above left=6mm and 10mm of ut] (fc) {FakeChroma\-Collection\\\scriptsize cosine in-memory};
\node[ext, below left=6mm and 10mm of ut] (hm) {HashingEmbedding\-Model\\\scriptsize deterministisch};
\draw[flowd] (fc) -- (ut);
\draw[flowd] (hm) -- (ut);

% real
\node[comp, right=42mm of ut] (val) {Retrieval-Validation\\\scriptsize 2 Tests};
\node[ext, above right=6mm and 10mm of val] (real) {all-MiniLM-L6-v2\\\scriptsize echtes Modell};
\node[store, below right=6mm and 10mm of val] (eph) {ChromaDB\\\scriptsize EphemeralClient};
\draw[flowd] (real) -- (val);
\draw[flowd] (eph) -- (val);

\node[lbl, below=1mm of ut] {kein Modell, keine I/O \textendash{} Coverage-Basis};
\node[lbl, below=1mm of val] {\texttt{skipif} wenn Modell/Deps fehlen};

\begin{scope}[on background layer]
\node[layer, fit=(fc)(hm)(ut), fill=soft,
  label={[font=\scriptsize\bfseries,color=accent]above:gemockt (CI-stabil)}]{};
\node[layer, fit=(real)(eph)(val), fill=accent3!8, draw=accent3,
  label={[font=\scriptsize\bfseries,color=accent3]above:echt (DoD-Nachweis)}]{};
\end{scope}
\end{tikzpicture}
\end{center}

\noindent Die \texttt{FakeChromaCollection} reproduziert ChromaDBs Cosine-Ranking,
sodass auch Retrieval-Logik ohne realen Store getestet wird. Die echte Validierung
läuft gegen einen \texttt{EphemeralClient} und das geladene Modell und überspringt
sauber, falls die schwere Stack-Installation fehlt.

% ===========================================================================
\section{Ergebnisse}
% ===========================================================================

\begin{minipage}[t]{0.5\linewidth}
\centering
\textbf{Test-Coverage je Modul}\\[2mm]
\begin{tikzpicture}
\begin{axis}[
  width=\linewidth, height=5.4cm,
  ybar, bar width=12pt,
  ymin=0, ymax=100,
  ylabel={\footnotesize Coverage \%},
  symbolic x coords={loader,chunker,embedder,vectordb},
  xtick=data, x tick label style={font=\footnotesize,rotate=18,anchor=east},
  ytick={0,20,40,60,80,100},
  nodes near coords, nodes near coords style={font=\tiny},
  enlarge x limits=0.18,
  axis lines=left,
  every axis plot/.append style={fill=accent,draw=accent},
]
\addplot coordinates {(loader,100)(chunker,100)(embedder,98)(vectordb,90)};
\draw[accent3,dashed,thick] (axis cs:loader,80) -- (axis cs:vectordb,80)
  node[above left,font=\tiny,color=accent3]{Ziel 80\%};
\end{axis}
\end{tikzpicture}
\end{minipage}\hfill
\begin{minipage}[t]{0.46\linewidth}
\centering
\textbf{Top-1 Cosine-Similarity (20 Queries)}\\[2mm]
\begin{tikzpicture}
\begin{axis}[
  width=\linewidth, height=5.4cm,
  ybar, bar width=20pt,
  ymin=0, ymax=1,
  ylabel={\footnotesize Similarity},
  symbolic x coords={min,Mittel,max},
  xtick=data, x tick label style={font=\footnotesize},
  nodes near coords, nodes near coords style={font=\tiny},
  enlarge x limits=0.3, axis lines=left,
  every axis plot/.append style={fill=accent2,draw=accent2},
]
\addplot coordinates {(min,0.479)(Mittel,0.663)(max,0.780)};
\draw[accent3,dashed,thick] (axis cs:min,0.25) -- (axis cs:max,0.25)
  node[above right,font=\tiny,color=accent3]{Threshold 0.25};
\end{axis}
\end{tikzpicture}
\end{minipage}

\vspace{4mm}
\noindent\textbf{Kernergebnis:} Auf dem 20-Dokument-Korpus erreicht das Retrieval eine
\textbf{Top-1-Hit-Rate von 100\,\%} \textendash{} jede Query findet als ähnlichsten
Treffer das thematisch korrekte Dokument, mit komfortablem Abstand zum Threshold.
Insgesamt laufen \textbf{32 Tests} grün; \texttt{flake8}, \texttt{mypy} und
\texttt{bandit} sind ohne Befund.

% ===========================================================================
\section{Betrieb \& Umgebung}
% ===========================================================================
Backend-Tests laufen über die Conda-Umgebung \texttt{Chat-Assistant-AI} (Python 3.11).
Drei umgebungsspezifische Stolpersteine wurden adressiert:

\begin{itemize}[leftmargin=1.4em]
  \item \textbf{Voller System-Datenträger:} pip- und HuggingFace-Temp/Cache werden auf
        einen beschreibbaren D:-Pfad umgeleitet (\texttt{TMP}, \texttt{HF\_HOME},
        \texttt{--cache-dir}).
  \item \textbf{Proxy-Zertifikat:} der Modell-Download nutzt ein kombiniertes
        CA-Bundle (certifi + Windows-Root-CAs) via \texttt{SSL\_CERT\_FILE}.
  \item \textbf{ChromaDB 1.x:} \texttt{delete\_collection} wirft \texttt{NotFoundError};
        \texttt{reset\_collection} fängt dies typsicher ab.
\end{itemize}

\paragraph{Neue Abhängigkeiten} (in \texttt{pyproject.toml}, Gruppen \texttt{backend} +
\texttt{development}): \texttt{sentence-transformers}, \texttt{pypdf},
\texttt{python-docx}, \texttt{langchain-text-splitters}.

\paragraph{Qualitäts-Kommandos}
\begin{lstlisting}[style=py,numbers=none]
flake8 --max-line-length=140 .
mypy --ignore-missing-imports backend/app/services/core/ingestion \
     backend/app/services/dependency/vectordb.py
bandit -r backend/app/services/core/ingestion
pytest tests/ragfeature --cov=backend/app/services/core/ingestion
\end{lstlisting}

\vfill
\begin{center}\small\color{codecom}
Ausblick: AP~2A (Gemini-Streaming-Chat) und AP~4 (vollständige LCEL-RAG-Chain mit
\texttt{retriever.py}) bauen direkt auf diesem Fundament auf.
\end{center}

\end{document}
